\documentclass[11pt,a4paper]{article}

\usepackage[utf8]{inputenc}
\usepackage[T1]{fontenc}
\usepackage[spanish,es-tabla]{babel}
\usepackage[a4paper,margin=2.2cm]{geometry}
\usepackage{xcolor}
\usepackage{graphicx}
\usepackage{booktabs}
\usepackage{longtable}
\usepackage{array}
\usepackage{tabularx}
\usepackage{enumitem}
\usepackage{titlesec}
\usepackage{fancyhdr}
\usepackage{hyperref}
\usepackage{listings}
\usepackage[font=small,labelfont=bf]{caption}

% Colores y estilo
\definecolor{mblue}{HTML}{1F4E79}
\definecolor{morange}{HTML}{C0392B}
\definecolor{mgreen}{HTML}{27AE60}
\definecolor{mgray}{HTML}{555555}
\definecolor{lightgray}{HTML}{F4F4F4}
\definecolor{codebg}{HTML}{F8F8F8}

\hypersetup{
  colorlinks=true,
  linkcolor=mblue,
  urlcolor=mblue,
  citecolor=mblue,
}

\titleformat{\section}{\Large\bfseries\color{mblue}}{\thesection}{1em}{}
\titleformat{\subsection}{\large\bfseries\color{mblue}}{\thesubsection}{1em}{}
\titleformat{\subsubsection}{\normalsize\bfseries\color{mgray}}{\thesubsubsection}{1em}{}

\pagestyle{fancy}
\fancyhf{}
\fancyhead[L]{\textcolor{mgray}{\small Reporte Técnico — Pipeline Xenium Lung Cancer}}
\fancyhead[R]{\textcolor{mgray}{\small INP — \today}}
\fancyfoot[C]{\textcolor{mgray}{\thepage}}
\renewcommand{\headrulewidth}{0.4pt}
\renewcommand{\headrule}{\hbox to\headwidth{\color{mblue}\leaders\hrule height \headrulewidth\hfill}}

\lstset{
  basicstyle=\ttfamily\footnotesize,
  backgroundcolor=\color{codebg},
  frame=single,
  rulecolor=\color{mgray},
  framesep=4pt,
  breaklines=true,
  breakatwhitespace=false,
  showstringspaces=false,
  numbers=none,
}

% Caja informativa
\newcommand{\infobox}[2]{%
  \par\noindent\fcolorbox{mblue}{lightgray}{\parbox{0.97\textwidth}{\textbf{\textcolor{mblue}{#1}}\\[2pt]#2}}\par
}

\begin{document}

% ════════════════════════════════════════════════════════════════════════════
% PORTADA
% ════════════════════════════════════════════════════════════════════════════
\thispagestyle{empty}
\vspace*{1.5cm}
\begin{center}
{\Huge\bfseries\color{mblue} Reporte Técnico\par}
\vspace{0.4cm}
{\Large\color{mgray} Pipeline de Análisis Xenium — Cáncer de Pulmón Humano\par}
\vspace{1cm}
\rule{0.7\textwidth}{0.6pt}\par
\vspace{0.6cm}
{\large
\begin{tabular}{rl}
\textbf{Autor:} & Miguel Ángel Díaz-Campos \\
\textbf{Institución:} & INP --- Instituto Nacional de Pediatría (México) \\
\textbf{Fecha:} & 29 de abril, 2026 \\
\textbf{Versión:} & v3 (anotación híbrida validada) \\
\end{tabular}
}
\vspace{1.2cm}

\begin{minipage}{0.85\textwidth}
\textbf{Resumen.} Este documento es el control técnico completo del pipeline de transcriptómica espacial Xenium aplicado a un dataset de cáncer de pulmón humano (268{,}034 células × 289 genes). Documenta: (i) la arquitectura del pipeline (Snakemake), (ii) la jerarquía de anotación celular L1/L2/L3 con 26 tipos granulares (versión 3 — anotación híbrida validada), (iii) la estrategia de validación de anotación (ARI vs subclustering computacional; CellTypist 5 modelos), (iv) las rutas absolutas de todos los archivos generados, (v) los formatos de datos y cómo cargarlos, y (vi) el estado de las herramientas instaladas.
\end{minipage}
\end{center}

\vspace{1cm}
\begin{center}
\begin{tabular}{|l|l|}
\hline
\rowcolor{lightgray}\textbf{Dataset} & Human lung cancer, Xenium v1, panel INP-Pulmón 2025 \\
\hline
\textbf{Células totales} & 268{,}034 \\
\textbf{Genes} & 289 \\
\textbf{Tipos celulares L1 (broad)} & 14 \\
\textbf{Tipos celulares L2 (granular)} & 26 \\
\textbf{Estados funcionales L3} & 7 \\
\textbf{Marcadores del panel} & 380 imm + 29 add-on lung \\
\textbf{Formato principal} & SpatialData / Zarr \\
\hline
\end{tabular}
\end{center}

\newpage
\tableofcontents
\newpage

% ════════════════════════════════════════════════════════════════════════════
\section{Estado General del Pipeline}

\subsection{Mapa de fases}

\begin{longtable}{|>{\bfseries}p{1cm}|p{4cm}|p{6.5cm}|p{2.5cm}|}
\hline
\rowcolor{lightgray}\textbf{Fase} & \textbf{Familia} & \textbf{Herramientas} & \textbf{Estado} \\
\hline
A & Snakemake skeleton + Phase 0 re-labeling v3 & Pipeline base + 26 tipos L2 (anotación híbrida validada) & \textcolor{mgreen}{✓ COMPLETO} \\
\hline
B & F8 Novae + F1 Banksy (dominios espaciales) & Foundation model + Banksy consensus & \textcolor{mgreen}{✓ COMPLETO} \\
\hline
C & F2 SVG Consensus & Hotspot + nnSVG + Moran's I & \textcolor{mgreen}{✓ COMPLETO} \\
\hline
D & F4 Niche-DE & NicheDE (Mason 2024) & \textcolor{morange}{✗ DESCARTADO} \\
\hline
E & F5 Spatial Pseudotime & Slingshot + tradeSeq + GAM & \textcolor{mgreen}{✓ COMPLETO 2026-04-30} \\
\hline
F & F3 CCC Validation & LIANA+ rank\_aggregate + Spacia & \textcolor{mgreen}{✓ COMPLETO 2026-04-30} \\
\hline
G & F6 Annotation x-val (CellTypist + ARI) & CellTypist 5 modelos + validación interna ARI & \textcolor{mgreen}{✓ COMPLETO 2026-04-29} \\
\hline
--- & F7 Co-expresión & ~~hdWGCNA / SCENIC~~ & \textcolor{morange}{✗ DESCARTADO} \\
\hline
\end{longtable}

\subsection{Decisiones críticas (descartes justificados)}

\begin{itemize}[leftmargin=*]
\item \textbf{F4 NicheDE descartado}: el panel de 289 genes es insuficiente; el método mismo lo confirmó (\textit{``Less than 1000 genes pass''}). Aplicable a Xenium Prime 5K (panel futuro).
\item \textbf{F7 Co-expresión descartada}: hdWGCNA / SCENIC requieren $\geq$3{,}000 genes para construir redes robustas.
\end{itemize}

\subsection{Anotación granular v3 — estrategia híbrida (2026-04-29)}

\infobox{¿Por qué v3 y no v2?}{
La anotación v2 usaba únicamente \texttt{sc.tl.score\_genes} argmax con los marcadores del panel. Al validar contra \texttt{cell\_type\_immune\_granular} (subclustering computacional Leiden sobre células inmunes, existente en \texttt{sdata.zarr}), el ARI fue solo \textbf{0.115} — clasificaciones erróneas: CD8\_T\_cell → cDC1 (31\%), Plasma\_cell → Unassigned (37\%). La v3 corrige esto con una estrategia híbrida de tres capas.
}

\textbf{Estrategia v3:}

\begin{enumerate}[leftmargin=*]
\item \textbf{Células inmunes (39.8\%, $n$=106,740)}: se usa directamente \texttt{cell\_type\_immune\_granular} (subclustering Leiden computacional, fuente más confiable). La granularidad CD8\_cytotoxic vs exhausted, B\_naive vs memory, cDC1 vs cDC2 vs pDC se resuelve por score argmax dentro de los candidatos del tipo.
\item \textbf{Clusters Leiden puros (21.5\%, $n$=57,650)}: Leiden~2 → Epithelial\_general (92\% pureza), Leiden~6 → Smooth\_muscle (88\%), Leiden~8 → Epithelial\_general (92\%). Asignación directa.
\item \textbf{Células no-inmunes restantes (38.7\%, $n$=103,644)}: \texttt{sc.tl.score\_genes} argmax limitado a tipos no-inmunes (AT1, Ciliated, Epithelial\_general, Smooth\_muscle, Pericyte, Endothelial, Tumor). Umbral $z>0.3$.
\end{enumerate}

\begin{center}
\begin{tabular}{|l|r|r|l|}
\hline
\rowcolor{lightgray}\textbf{Comparación} & \textbf{ARI v2} & \textbf{ARI v3} & \textbf{Mejora} \\
\hline
L2 vs immune\_granular & 0.115 & \textbf{0.742} & +543\% \\
L2 vs Leiden global & 0.263 & 0.293 & +11\% \\
L2 vs leiden\_immune & 0.116 & 0.194 & +67\% \\
\hline
\end{tabular}
\end{center}

\infobox{Nota sobre CellTypist}{
Se corrieron 5 modelos CellTypist (\texttt{Human\_Lung\_Atlas}, \texttt{Human\_PF\_Lung}, \texttt{Human\_IPF\_Lung}, \texttt{Immune\_All\_High}, \texttt{Immune\_All\_Low}) para validación cruzada. Todos colapsaron a 4--6 categorías amplias porque el panel de 289 genes tiene cobertura insuficiente para clasificadores entrenados con transcriptoma completo ($\sim$20{,}000 genes). \textbf{CellTypist no es un validador apropiado para paneles Xenium v1 dirigidos.} La métrica de validación correcta es el ARI interno vs subclustering computacional.
}

% ════════════════════════════════════════════════════════════════════════════
\newpage
\section{Datos: Formatos, Rutas y Acceso}

\subsection{Objeto principal: SpatialData (Zarr)}

El dataset principal vive en formato \textbf{SpatialData} (estándar scverse 2024), persistido como Zarr (directorio jerárquico). Contiene tablas, imágenes, formas (cell boundaries), puntos (transcripts) y metadatos.

\begin{center}
\begin{tabularx}{\textwidth}{|l|X|}
\hline
\rowcolor{lightgray}\textbf{Tipo} & \textbf{Ruta absoluta} \\
\hline
\textbf{ACTIVO (v3)} & \texttt{\small /home/mdiaz/Documents/Xenium\_project/proyecto\_demo\_xenium/human\_lung\_cancer/results/sdata.zarr} \\
\hline
\textcolor{morange}{\textbf{BACKUP (no usar)}} & \texttt{\small /home/mdiaz/Documents/Xenium\_project/proyecto\_demo\_xenium/human\_lung\_cancer/results/sdata\_DESACTUALIZADO\_NO\_USAR.zarr} \\
\hline
\end{tabularx}
\end{center}

\subsubsection{Cómo cargar el objeto}

\begin{lstlisting}[language=Python]
import spatialdata as sd
sdata = sd.read_zarr(
    "/home/mdiaz/Documents/Xenium_project/proyecto_demo_xenium/"
    "human_lung_cancer/results/sdata.zarr"
)
adata = sdata.tables["table"]   # AnnData con 268,034 cells x 289 genes
print(adata.obs.columns)         # 66 columnas (incluye L1, L2, L3, scores)
print(adata.obs["cell_type_L2"].value_counts())
\end{lstlisting}

\subsection{Estructura interna del SpatialData}

\begin{longtable}{|>{\bfseries}p{3.5cm}|p{4cm}|p{6cm}|}
\hline
\rowcolor{lightgray}\textbf{Componente} & \textbf{Clave} & \textbf{Contenido} \\
\hline
\texttt{tables} & \texttt{table} & AnnData (268{,}034 × 289), 66 cols obs, 1 layer (counts), obsm (PCA, UMAP, spatial) \\
\hline
\texttt{images} & \texttt{he\_image}, \texttt{morphology\_focus} & H\&E + DAPI multiscale (OME-Zarr) \\
\hline
\texttt{shapes} & \texttt{cell\_boundaries}, \texttt{cell\_circles}, \texttt{nucleus\_boundaries} & Polígonos de segmentación (xoa) \\
\hline
\texttt{points} & \texttt{transcripts} & Coordenadas espaciales de cada transcrito detectado \\
\hline
\end{longtable}

\subsection{Columnas de anotación (post-v3)}

\begin{longtable}{|>{\ttfamily}l|p{2.5cm}|p{8cm}|}
\hline
\rowcolor{lightgray}\textbf{Columna} & \textbf{Categorías} & \textbf{Descripción} \\
\hline
cell\_type\_L1 & 14 & Broad: Epithelial\_alveolar/airway/general, Stromal, T\_cell, Macrophage, Endothelial, Dendritic\_cell, Tumor, Plasma\_cell, Monocyte, B\_cell, NK\_cell, Unassigned \\
\hline
cell\_type\_L2 & 26 & Granular (ver Sección 3) — \textbf{canónica v3} \\
\hline
cell\_type\_L3 & 7 & Estados: Antigen\_presenting, Cytotoxic\_effector, Exhaustion, Immune\_regulatory, Pro\_inflammatory, Proliferating, Steady\_state \\
\hline
cell\_type & 14 & Alias de L1 (compatibilidad) \\
\hline
cell\_type\_fine & 26 & Alias de L2 (compatibilidad) \\
\hline
cell\_type\_L1\_v3 / \_L2\_v3 / \_L3\_v3 & — & Columnas v3 originales (conservadas para trazabilidad) \\
\hline
annotation\_source\_v3 & categórica & Fuente: \texttt{immune\_granular} / \texttt{leiden\_X\_direct} / \texttt{score\_nonimmune} \\
\hline
cell\_type\_score & numeric & z-score top-1 (sc.tl.score\_genes, para células con fuente score) \\
\hline
cell\_type\_confidence & 4 niveles & LOW / MEDIUM / HIGH / VERY\_HIGH (gap top1 - top2) \\
\hline
score\_<celltype> & 26 cols & Score individual por tipo L2 \\
\hline
\end{longtable}

\subsection{Layers, OBSM, UNS}

\begin{itemize}[leftmargin=*]
\item \texttt{adata.layers["counts"]} --- matriz dispersa CSR de counts crudos (268{,}034 × 289)
\item \texttt{adata.X} --- normalizado + log1p (para PCA, scoring, clustering)
\item \texttt{adata.obsm["spatial"]} --- coordenadas (x, y) en \(\mu\)m
\item \texttt{adata.obsm["X\_pca"]} --- 30 componentes principales
\item \texttt{adata.obsm["X\_umap"]} --- 2 componentes UMAP
\item \texttt{adata.uns["moranI"]} --- Moran's I por gene (Phase 1A)
\item \texttt{adata.uns["liana\_res"]} --- resultados LIANA+ (sobrescritos al re-correr)
\end{itemize}

% ════════════════════════════════════════════════════════════════════════════
\newpage
\section{Anotación Celular Granular (v3 — híbrida validada)}

\subsection{Marcadores del panel INP-Pulmón 2025}

\textbf{Fuente:} \texttt{\small /home/mdiaz/Documents/Xenium\_project/my\_xenium\_panel\_markers/Propuesta panel Pulmón 2025.xlsx - Hoja1.csv}

\begin{itemize}[leftmargin=*]
\item \textbf{Inmuno panel:} 380 genes con anotación de tipo celular (T cell, B cell, NK, Macrophage, DC, etc.)
\item \textbf{Add-on lung:} 29 genes con tejido (Lung / Fibrotic Lung) y tipo (AT1, Ciliated, Club, Fibroblast, Basaloid)
\end{itemize}

\subsection{Diccionario de scoring (extracto)}

El diccionario completo está en \texttt{F0\_reannotation\_v3.py} variable \texttt{MARKERS}. Resumen:

\begin{longtable}{|>{\ttfamily}l|p{8cm}|}
\hline
\rowcolor{lightgray}\textbf{Tipo L2} & \textbf{Marcadores usados} \\
\hline
CD4\_T\_helper & CD4, IL7R, CCR7, TCF7, LEF1, CD3D, CD3E, CD3G \\
CD8\_T\_cytotoxic & CD8A, CD8B, GZMA, GZMB, GZMK, PRF1, NKG7, CD3D \\
CD8\_T\_exhausted & CD8A, PDCD1, LAG3, HAVCR2, TIGIT, TOX, CTLA4 \\
Treg & FOXP3, IL2RA, CTLA4, IKZF2, TIGIT, CD4 \\
B\_naive & MS4A1, CD19, CD79A, CD79B, IGHD, TCL1A \\
B\_memory & MS4A1, CD19, CD27, CD79A, CD79B \\
Plasma & MZB1, JCHAIN, XBP1, CD27, TNFRSF17, PRDM1 \\
Plasmablast & MKI67, MZB1, JCHAIN, XBP1 \\
NK\_cytotoxic & NCAM1, FCGR3A, KLRD1, NKG7, GNLY, PRF1, GZMB, KLRF1 \\
Macrophage\_M1 & CD68, CD86, TNF, IL1B, NOS2, CXCL9, CXCL10, AIF1 \\
Macrophage\_M2 & CD68, CD163, MRC1, MSR1, VSIG4, MARCO, AIF1, CD14 \\
Monocyte\_classical & CD14, S100A8, S100A9, FCN1, VCAN, CSF3R \\
cDC1 & CLEC9A, XCR1, BATF3, IRF8, CD8A \\
cDC2 & CD1C, FCER1A, CLEC10A, CD1A \\
DC\_mature & LAMP3, CCR7, CCL19, CCL22, FSCN1 \\
pDC & LILRA4, IRF7, IRF8, CLEC4C, TCF4 \\
AT1 & AGER, CAVIN1, PDPN, RTKN2 \\
Ciliated & FOXJ1, AGR3, C20orf85, C1orf194, C6orf118, CCDC39, CCDC78, CFAP53 \\
Smooth\_muscle & ACTA2, MYH11, TAGLN, DES, CNN1 \\
Pericyte & RGS5, PDGFRB, MCAM, ACTA2, NOTCH3 \\
Endothelial\_blood & PECAM1, VWF, CDH5, CLDN5, ADGRL4, ACKR1 \\
Endothelial\_lymphatic & PECAM1, PROX1, PDPN, LYVE1, TFF3 \\
Tumor\_proliferating & TOP2A, MKI67, CCNB1, CDK1, BIRC5 \\
Tumor\_resting & EPCAM, KRT8, KRT18, MUC1, TSPAN8, WFS1, AGR3 \\
\hline
\end{longtable}

\textbf{Nota:} Marcadores diseñados para 35 tipos; \textbf{26 pasaron} el filtro de $\geq$2 marcadores presentes en el panel (los excluidos por baja cobertura: T\_resident\_memory, NK\_resting, Mast\_cell, Neutrophil, Goblet, Basal, Basaloid\_aberrant, AT2, Fibroblast, Club).

\subsection{Algoritmo de asignación}

\begin{enumerate}[leftmargin=*]
\item Para cada tipo $T_i \in \{T_1, ..., T_{26}\}$ con marcadores $G_i$:
$$ \text{score}_{T_i}(c) = \frac{1}{|G_i|}\sum_{g\in G_i} x_{c,g} - \frac{1}{|R_i|}\sum_{g\in R_i} x_{c,g} $$
donde $R_i$ es un conjunto random de genes con expresión similar (\texttt{ctrl\_size=50}).
\item Asignación: $L2(c) = \arg\max_i \text{score}_{T_i}(c)$ si $\max_i > 0.5$, else \texttt{Unassigned}.
\item Confianza: $\text{conf}(c) = \text{score}_{(1)} - \text{score}_{(2)}$ \\
LOW $<0.1$ · MEDIUM $[0.1, 0.3)$ · HIGH $[0.3, 0.6)$ · VERY\_HIGH $\geq 0.6$.
\end{enumerate}

\subsection{Distribución de tipos celulares (v3)}

\subsubsection{L1 (14 categorías)}
Epithelial\_general 22.1\% · T\_cell 20.6\% · Stromal 12.0\% · Endothelial 9.6\% · Macrophage 8.5\% · Epithelial\_alveolar 5.1\% · Unassigned 4.8\% · Tumor 4.4\% · B\_cell 3.9\% · Monocyte 3.1\% · Epithelial\_airway 2.3\% · NK\_cell 2.0\% · Dendritic\_cell 1.3\% · Plasma\_cell 0.3\%.

\subsubsection{L2 (26 tipos — distribución completa v3)}
\begin{tabular}{lrr}
\rowcolor{lightgray}\textbf{Tipo L2} & \textbf{N células} & \textbf{\%} \\
Epithelial\_general & 59{,}319 & 22.1\% \\
CD4\_T\_helper & 31{,}378 & 11.7\% \\
Endothelial\_blood & 21{,}225 & 7.9\% \\
Smooth\_muscle & 19{,}506 & 7.3\% \\
Macrophage\_M2 & 18{,}873 & 7.0\% \\
Treg & 16{,}342 & 6.1\% \\
AT1 & 13{,}572 & 5.1\% \\
Unassigned & 12{,}782 & 4.8\% \\
Pericyte & 12{,}587 & 4.7\% \\
Tumor\_resting & 8{,}819 & 3.3\% \\
Monocyte\_classical & 8{,}422 & 3.1\% \\
CD8\_T\_cytotoxic & 6{,}284 & 2.3\% \\
Ciliated & 6{,}035 & 2.3\% \\
B\_naive & 6{,}018 & 2.2\% \\
NK\_cytotoxic & 5{,}439 & 2.0\% \\
Endothelial\_lymphatic & 4{,}549 & 1.7\% \\
B\_memory & 4{,}399 & 1.6\% \\
Macrophage\_M1 & 4{,}029 & 1.5\% \\
Tumor\_proliferating & 2{,}900 & 1.1\% \\
cDC1 & 1{,}807 & 0.7\% \\
cDC2 & 1{,}290 & 0.5\% \\
CD8\_T\_exhausted & 1{,}228 & 0.5\% \\
Plasmablast & 431 & 0.2\% \\
pDC & 395 & 0.1\% \\
Plasma & 307 & 0.1\% \\
DC\_mature & 98 & 0.04\% \\
\end{tabular}

\vspace{4pt}
\textbf{Total biológico (sin Unassigned): 25 tipos, 255{,}252 células (95.2\%)}

\textbf{Fuente de anotación:} immune\_granular 39.8\% · score\_nonimmune 38.7\% · leiden\_directo 21.5\%

% ════════════════════════════════════════════════════════════════════════════
\newpage
\section{Resultados Generados (rutas absolutas)}

\subsection{Anotación v3 (canónica)}
\textbf{Directorio:} \texttt{\small /home/mdiaz/.../human\_lung\_cancer/results/02\_biology/reannotation\_v3/}

\begin{tabular}{|>{\ttfamily}l|p{6.5cm}|}
\hline
\rowcolor{lightgray}\textbf{Archivo} & \textbf{Contenido} \\
\hline
annotation\_v3\_per\_cell.csv & Por célula: L2\_v3, L1\_v3, annotation\_source\_v3 \\
distribution\_v3.csv & Distribución: L2, L1, n\_cells, pct \\
annotation\_sources\_v3.csv & Recuento por fuente \\
annotation\_v3\_summary.json & Resumen + modelos de referencia utilizados \\
\hline
\end{tabular}

\subsection{Validación de anotación (F6)}
\textbf{Directorio:} \texttt{\small /home/mdiaz/.../human\_lung\_cancer/results/02\_biology/F6\_celltypist\_validation/}

\begin{tabular}{|>{\ttfamily}l|p{6.5cm}|}
\hline
\rowcolor{lightgray}\textbf{Archivo} & \textbf{Contenido} \\
\hline
concordance\_summary.json & ARI/NMI: L2 vs Leiden, vs immune\_granular, vs leiden\_immune \\
internal\_L2\_vs\_leiden.csv & Composición L2 por cluster Leiden \\
internal\_L2\_vs\_immune\_granular.csv & L2 vs subclustering inmune (tabla confusión) \\
heatmap\_L2\_vs\_leiden.png & Heatmap composición \\
heatmap\_L2\_vs\_immune\_granular.png & Heatmap concordancia \\
confusion\_Human\_Lung\_Atlas.pkl.csv & CellTypist (documentación — modelo insuf.) \\
\hline
\end{tabular}

\subsection{DGEA L2 v3 (Wilcoxon 1 vs Rest)}
\textbf{Directorio:} \texttt{\small /home/mdiaz/.../human\_lung\_cancer/results/02\_biology/dgea\_L2\_v3/}

\begin{itemize}[leftmargin=*]
\item \textbf{24} CSVs (24 tipos con $\geq$100 células): \texttt{dgea\_L2\_<tipo>.csv}
\item \texttt{dgea\_L2\_top50\_all.csv} --- top-50 por tipo concatenado
\item \texttt{dgea\_L2\_dotplot\_top5.png} --- dotplot top-5 marcadores por tipo
\item \texttt{dgea\_L2\_summary.json} --- 255{,}154 células usadas, 24 grupos retenidos
\end{itemize}

\subsection{CCC LIANA+ L2 v3 (rank\_aggregate)}
\textbf{Directorio:} \texttt{\small /home/mdiaz/.../human\_lung\_cancer/results/02\_biology/ccc\_liana\_L2\_v3/}

\begin{itemize}[leftmargin=*]
\item \texttt{liana\_L2\_all\_interactions.csv} --- 2{,}631 interacciones totales (sin self)
\item \texttt{liana\_L2\_significant.csv} --- \textbf{1{,}565 significativas} (filtro \texttt{magnitude\_rank})
\item \texttt{liana\_L2\_top50.csv} --- top-50 hits por magnitud
\item \texttt{liana\_L2\_celltype\_pair\_counts.csv} --- matriz $25 \times 25$ de \# pares L-R por par L2
\item \texttt{liana\_L2\_heatmap.png} --- visualización del crosstalk
\item \texttt{liana\_L2\_summary.json} --- ejecución 50.4s, n\_perms=1000
\end{itemize}

\subsection{Phase B (dominios espaciales)}
\textbf{Directorio:} \texttt{\small /home/mdiaz/.../human\_lung\_cancer/results/03\_phase3\_spatial/}

\begin{itemize}[leftmargin=*]
\item \texttt{F1\_spatial\_domains/} --- Banksy (14 dominios) + consenso F1+F8
\item \texttt{F8\_novae/} --- Embeddings 64d Novae + dominios L5/L10/L20 (multi-resolución)
\end{itemize}

\subsection{Phase C (SVG consensus)}
\textbf{Directorio:} \texttt{\small /home/mdiaz/.../human\_lung\_cancer/results/03\_phase3\_spatial/F2\_svg\_consensus/}

\begin{itemize}[leftmargin=*]
\item \texttt{hotspot\_svg\_scores.csv} (289 genes), \texttt{nnsvg\_svg\_scores.csv} (10k subsample)
\item \texttt{svg\_consensus\_table.csv}: 288/289 genes en $\geq$2 métodos
\item \texttt{svg\_venn.png}, \texttt{svg\_consensus\_heatmap.png}
\end{itemize}

\subsection{Phase E (Spatial Pseudotime)}
\textbf{Completado:} 2026-04-30 | \textbf{Entorno:} \texttt{xenium\_R\_analysis}

\textbf{Método:} Slingshot + tradeSeq sobre embeddings Novae 64d (29,991 células estratificadas del TME).

\begin{itemize}[leftmargin=*]
\item \textbf{8 linajes} identificados desde \texttt{Tumor\_proliferating} como punto de origen
\item \textbf{Linaje 1:} 17,612 células, rango pseudotime 0 $\to$ 5.28
\item \textbf{tradeSeq GAMs:} todas las 20 dimensiones Novae con $p=0$ --- señal pseudotemporal robusta en todo el espacio latente
\item \textbf{Figuras generadas:}
  \begin{itemize}
  \item \texttt{F5\_umap\_celltypes.png} --- UMAP coloreado por tipo celular L1
  \item \texttt{F5\_umap\_pseudotime\_L1.png} --- UMAP coloreado por pseudotime
  \item \texttt{F5\_spatial\_pseudotime.png} --- mapa espacial de pseudotime
  \item \texttt{F5\_pseudotime\_vs\_spatial\_gradient.png} --- gradiente vs distancia al tumor
  \end{itemize}
\end{itemize}

\textbf{Directorio:} \texttt{\small human\_lung\_cancer/results/figures/phase5\_pseudotime/} \\
\textbf{Script:} \texttt{pipeline/scripts/analysis/F5\_pseudotime\_slingshot.R}

\subsection{Phase F (CCC Validation --- LIANA+ × Spacia)}
\textbf{Completado:} 2026-04-30 | \textbf{Entorno:} \texttt{spacia} (Bayesian MIL)

\textbf{Método:} Spacia (Zhu 2024, \textit{Nature Methods}) valida si las interacciones L-R de LIANA+ tienen señal espacial real. Para cada par (sender, receiver, ligando, receptor), ajusta un modelo Bayesian Multiple Instance Learning: el ligando en el sender predice la expresión del receptor en células receptoras dentro de 30 µm.

\textbf{Criterio de validación:} \texttt{Pathway\_betas.csv} $\to$ pval\_adj $< 0.05$.

\textbf{Resultados globales:}
\begin{itemize}[leftmargin=*]
\item \textbf{30 interacciones} seleccionadas (top LIANA+ con ambos genes en panel)
\item \textbf{25 jobs completados} (5 FAILs por tipos raros: Plasma, pDC$\times$1, Plasmablast, DC\_mature --- insuficientes bags MIL con $<500$ células sender)
\item \textbf{22/25 validadas espacialmente (88\%)}
\end{itemize}

\textbf{Hallazgos principales:}

\begin{longtable}{|p{3.5cm}|p{2.2cm}|p{2.5cm}|p{2.5cm}|r|}
\hline
\rowcolor{lightgray}\textbf{Sender} & \textbf{Receiver} & \textbf{Ligando} & \textbf{Receptor} & \textbf{pval\_adj} \\
\hline
AT1 & Tumor\_resting & ADAM17 & MUC1 & $1.1\times10^{-71}$ \\
\hline
Tumor\_resting & AT1 & LTF & AGER & $9.4\times10^{-66}$ \\
\hline
Monocyte\_classical & Tumor\_resting & ADAM17 & MUC1 & $8.4\times10^{-55}$ \\
\hline
Pericyte & Tumor\_resting & ADAM17 & MUC1 & $3.5\times10^{-48}$ \\
\hline
Macrophage\_M1 & Epithelial\_general & ADAM17 & MUC1 & $5.9\times10^{-43}$ \\
\hline
Epithelial\_general & Endothelial\_lymphatic & CDH1 & EGFR & $6.6\times10^{-42}$ \\
\hline
cDC2 & AT1 & S100B & AGER & $4.0\times10^{-14}$ \\
\hline
Ciliated & AT1 & LTF & AGER & $1.2\times10^{-4}$ \\
\hline
\end{longtable}

\textbf{Ejes de señalización validados:}
\begin{itemize}[leftmargin=*]
\item \textbf{ADAM17 $\to$ MUC1 (hub central):} 7+ senders (AT1, M1, M2, Monocyte\_classical, Pericyte, cDC1, cDC2, Epithelial\_general) $\to$ Tumor\_resting. Independiente del estado de polarización macrofágica (M1 y M2 ambos validan). Relevante para remodelación del microambiente tumoral.
\item \textbf{LTF $\to$ AGER:} Tumor\_resting y Ciliated $\to$ AT1. Señalización tumor$\to$epitelio alveolar (invasión del nicho).
\item \textbf{S100B $\to$ AGER:} cDC2, cDC1, Tumor\_proliferating $\to$ AT1. Señalización inflamatoria DC-mediada.
\item \textbf{CDH1 $\to$ EGFR:} Epithelial\_general como sender hub $\to$ AT1, Ciliated, Plasmablast, Endothelial\_lymphatic. Eje E-cadherina/EGFR transversal.
\end{itemize}

\textbf{Directorios:}
\begin{itemize}[leftmargin=*]
\item Jobs: \texttt{\small human\_lung\_cancer/results/02\_biology/phase\_f\_spacia/jobs/}
\item Resultados: \texttt{\small phase\_f\_spacia/F3\_spacia\_results.csv}
\item Figuras: \texttt{\small results/figures/phase\_f\_spacia/F3\_spacia\_validation\_bubble.png}
\end{itemize}

\textbf{Scripts:} \texttt{F3\_spacia\_ccc\_validation.py}, \texttt{F3b\_run\_all\_spacia\_jobs.sh}, \texttt{F3c\_aggregate\_spacia\_results.py}

\subsection{Master Figure (5 paneles integrados)}
\textbf{Completado:} 2026-04-30 | \textbf{Script:} \texttt{F\_master\_figure.py}

\begin{itemize}[leftmargin=*]
\item \textbf{Panel A:} Mapa espacial de dominios Novae + Banksy (268,034 células; consenso 35.4\% resaltado)
\item \textbf{Panel B:} SVG consensus heatmap (Hotspot + nnSVG + Moran's I; 288/289 genes)
\item \textbf{Panel C:} Validación CCC LIANA+ × Spacia (22/25 interacciones, bar chart $-\log_{10}$ pval\_adj)
\item \textbf{Panel D:} UMAP pseudotime Slingshot (8 linajes, 29,991 células)
\item \textbf{Panel E:} Heatmap annotation x-val ARI = 0.742
\end{itemize}

\textbf{Archivos:} \texttt{\small results/figures/master\_figure/master\_figure\_v1.png} (300 DPI) + \texttt{.pdf}

% ════════════════════════════════════════════════════════════════════════════
\newpage
\section{Scripts del Pipeline}

\subsection{Scripts activos (válidos post-v3)}

\begin{longtable}{|>{\ttfamily}p{6.5cm}|p{2.5cm}|p{4.5cm}|}
\hline
\rowcolor{lightgray}\textbf{Script} & \textbf{Entorno} & \textbf{Función} \\
\hline
F0\_reannotation\_v3.py & xenium\_pipeline & \textbf{(CANÓNICO)} Anotación híbrida: immune\_granular + Leiden + score \\
\hline
F6\_celltypist\_validation.py & xenium\_pipeline & CellTypist 5 modelos (documentación, no usar como validador) \\
\hline
F6b\_internal\_annotation\_validation.py & xenium\_pipeline & ARI/NMI interno: L2 vs Leiden vs immune\_granular \\
\hline
F0b\_dgea\_granular\_v2.py & xenium\_pipeline & DGEA Wilcoxon sobre cell\_type\_L2 \\
\hline
F0c\_ccc\_liana\_granular\_v2.py & xenium\_pipeline & LIANA+ rank\_aggregate sobre L2 \\
\hline
F1\_export\_for\_banksy.py + F1\_banksy\_domains.R & ambos & Banksy spatial domains \\
\hline
F1\_novae\_domains.py & xenium\_pipeline & Novae multi-resolución (L5/L10/L20) \\
\hline
F1\_domains\_consensus.py & xenium\_pipeline & Consenso Banksy + Novae (Hungarian) \\
\hline
F8\_novae\_embeddings.py & xenium\_pipeline & Foundation model embeddings 64d \\
\hline
F2\_export\_for\_nnsvg.py + F2\_nnsvg.R & ambos & nnSVG en 10k subsample \\
\hline
F2\_hotspot\_svg.py & xenium\_pipeline & Hotspot SVG en 268k completo \\
\hline
F2\_svg\_consensus.py & xenium\_pipeline & Integración 3 métodos SVG \\
\hline
F5\_pseudotime\_slingshot.R & xenium\_R\_analysis & Slingshot + tradeSeq + GAM (8 linajes, Novae 64d) \\
\hline
F3\_spacia\_ccc\_validation.py & spacia & Pipeline Spacia Bayesian MIL (30 interacciones LIANA+) \\
\hline
F3b\_run\_all\_spacia\_jobs.sh & spacia & Runner secuencial 30 jobs Spacia \\
\hline
F3c\_aggregate\_spacia\_results.py & xenium\_pipeline & Agrega resultados Spacia + figuras validación \\
\hline
F\_master\_figure.py & xenium\_pipeline & Master Figure 5 paneles (PNG 300 DPI + PDF) \\
\hline
\end{longtable}

\subsection{Pipeline Snakemake}

\textbf{Snakefile principal (preprocesamiento, no tocar):} \\
\texttt{\small /home/mdiaz/Documents/Xenium\_project/proyecto\_demo\_xenium/pipeline/Snakefile}

\textbf{Snakefile de análisis (DAG F1-F8):} \\
\texttt{\small /home/mdiaz/Documents/Xenium\_project/proyecto\_demo\_xenium/pipeline/Snakefile\_analysis}

\textbf{Reglas por familia:} \texttt{\small pipeline/analysis\_rules/F\{0..8\}\_*.smk}

\subsection{Comando de ejecución}
\begin{lstlisting}[language=bash]
cd /home/mdiaz/Documents/Xenium_project/proyecto_demo_xenium/pipeline
conda activate xenium_pipeline
snakemake -s Snakefile_analysis --configfile config/config_lung.yaml \
          --cores 8 --use-conda
\end{lstlisting}

% ════════════════════════════════════════════════════════════════════════════
\newpage
\section{Herramientas Instaladas}

\subsection{Entorno \texttt{xenium\_pipeline} (Python)}

\begin{tabular}{|l|l||l|l|}
\hline
\rowcolor{lightgray}\textbf{Herramienta} & \textbf{Versión} & \textbf{Herramienta} & \textbf{Versión} \\
\hline
scanpy & 1.11.5 & matplotlib & 3.10.8 \\
spatialdata & 0.7.2 & numpy & 2.4.4 \\
anndata & 0.12.10 & pandas & 2.3.3 \\
squidpy & 1.8.1 & scipy & 1.17.1 \\
liana & 1.7.1 & sklearn & 1.8.0 \\
hotspot (hotspotsc) & 1.1.3 & numba & 0.65.0 \\
novae & 1.0.5 & igraph & 1.0.0 \\
celltypist & 1.7.1 & sklearn & 1.8.0 \\
\hline
\end{tabular}

\subsection{Entorno \texttt{xenium\_R\_analysis} (R)}

\begin{tabular}{|l|l||l|l|}
\hline
\rowcolor{lightgray}\textbf{Herramienta} & \textbf{Versión} & \textbf{Herramienta} & \textbf{Versión} \\
\hline
Banksy & 1.6.0 & Matrix & 1.7.5 \\
nnSVG & 1.14.0 & Rfast & 2.1.5.2 \\
slingshot & 2.18.0 & BiocParallel & 1.44.0 \\
tradeSeq & 1.24.0 & doMC / doParallel & 1.3.8 / 1.0.17 \\
SingleR & 2.12.0 & foreach & 1.5.2 \\
celldex & 1.20.0 & scuttle & 1.20.0 \\
SpatialExperiment & 1.20.0 & bluster & 1.20.0 \\
scran & 1.38.1 & metapod & 1.18.0 \\
scater & 1.38.1 & RcppML & 0.3.7.1 \\
\hline
\end{tabular}

\subsection{Entorno \texttt{spacia} (Python 3.8 + R)}

Aislado para Spacia (Bayesian MIL CCC validation, Zhu 2024). Sólo se activa para Phase F.

% ════════════════════════════════════════════════════════════════════════════
\newpage
\section{Cómo Cargar y Explorar los Datos}

\subsection{Inspección rápida (Python)}

\begin{lstlisting}[language=Python]
# /home/mdiaz/.conda/envs/xenium_pipeline/bin/python
import spatialdata as sd
sdata = sd.read_zarr("/home/mdiaz/Documents/Xenium_project/proyecto_demo_xenium/"
                      "human_lung_cancer/results/sdata.zarr")
adata = sdata.tables["table"]

# 1) Tipos celulares granulares
print(adata.obs["cell_type_L2"].value_counts())
# 2) Confianza
print(adata.obs["cell_type_confidence"].value_counts())
# 3) Coordenadas espaciales
import numpy as np
coords = adata.obsm["spatial"]
print(f"x range: {coords[:,0].min():.1f}-{coords[:,0].max():.1f}")
# 4) UMAP
import scanpy as sc
sc.pl.umap(adata, color="cell_type_L2", legend_loc="on data")
# 5) Score de un tipo
sc.pl.umap(adata, color=["score_Treg", "score_Macrophage_M2"])
\end{lstlisting}

\subsection{Cargar resultados generados}

\begin{lstlisting}[language=Python]
import pandas as pd
# DGEA top markers Macrophage_M2
m2 = pd.read_csv("/home/mdiaz/Documents/Xenium_project/proyecto_demo_xenium/"
                 "human_lung_cancer/results/02_biology/dgea_L2_v2/dgea_L2_Macrophage_M2.csv")
print(m2.head(10)[["names", "scores", "logfoldchanges", "pvals_adj"]])

# CCC LIANA+ top hits
ccc = pd.read_csv("/home/mdiaz/Documents/Xenium_project/proyecto_demo_xenium/"
                  "human_lung_cancer/results/02_biology/ccc_liana_L2_v2/liana_L2_top50.csv")
print(ccc[["source","target","ligand_complex","receptor_complex","magnitude_rank"]].head(15))
\end{lstlisting}

\subsection{Acceso desde R (vía SpatialExperiment)}

\begin{lstlisting}[language=R]
# /home/mdiaz/.conda/envs/xenium_R_analysis/bin/Rscript
library(SpatialExperiment); library(zellkonverter)
spe <- readH5AD("/path/to/exported.h5ad")  # exportar previamente con anndata.write_h5ad
\end{lstlisting}

% ════════════════════════════════════════════════════════════════════════════
\newpage
\section{Interpretación Biológica (resumen)}

\subsection{Microambiente tumoral pulmonar revelado}

\begin{itemize}[leftmargin=*]
\item \textbf{Eje supresivo dominante (Treg + M2 macrophage):} Treg representa 4.2\% y M2 4.3\% del total. Las 1{,}710 interacciones LIANA+ confirman que estos dos tipos comparten señalización inmunorregulatoria (CTLA4, TIGIT, IL10).
\item \textbf{Componente epitelial heterogéneo:} se diferenciaron AT1 (5.9\%, parénquima alveolar normal residual), Ciliated (2.2\%, vía aérea), Tumor\_resting (5.5\%) y Tumor\_proliferating (1.0\%). El gradiente tumoral está separado de la arquitectura pulmonar normal.
\item \textbf{Estroma + vasculatura:} Stromal (Smooth\_muscle + Pericyte + Fibroblast) = 13.5\%, Endothelial (blood + lymphatic) = 7.4\%. Microvasculatura activa en el tejido.
\item \textbf{Compartimento adaptativo activo:} CD4\_T\_helper (7.4\%) > Treg (4.2\%) > CD8\_T\_cytotoxic (1.0\%) > CD8\_T\_exhausted (0.4\%). Predominio CD4 sobre CD8 sugiere respuesta adaptativa Th2/inmunoregulada en lugar de citotóxica.
\item \textbf{Compartimento mieloide diverso:} M2 (4.3\%) > M1 (3.4\%); cDC1 (3.5\%) > DC\_mature (2.2\%) > pDC (0.7\%) > cDC2 (0.4\%). Maduración DC presente pero subdominante.
\end{itemize}

\subsection{Implicación terapéutica}

El microambiente está dominado por mecanismos mieloide-supresivos (M2 + Treg + cDC inmadura) más que por exhaustion clásica de CD8. Sugiere mayor potencial para terapias dirigidas a re-polarización mieloide (anti-CSF1R, anti-CD163) que para checkpoint blockade clásico anti-PD1/anti-PD-L1.

% ════════════════════════════════════════════════════════════════════════════
\newpage
\section{Estado Final del Pipeline (2026-04-30)}

\textbf{El pipeline multi-herramienta está COMPLETO.} Todas las fases planificadas han sido ejecutadas o descartadas con justificación.

\begin{itemize}[leftmargin=*]
\item \textbf{Completadas:} A (anotación v3), B (dominios espaciales), C (SVG consensus), E (pseudotime), F (CCC validation), G (annotation x-val) + Master Figure
\item \textbf{Descartadas justificadamente:} D (NicheDE, 289 genes insuficientes), F7 (hdWGCNA/SCENIC, misma razón)
\item \textbf{Pendiente único:} Preparar manuscrito / submisión
\end{itemize}

\subsection*{Métricas clave del pipeline}

\begin{tabular}{|l|r|}
\hline
\rowcolor{lightgray}\textbf{Métrica} & \textbf{Valor} \\
\hline
Células analizadas & 268,034 \\
\hline
Genes en panel & 289 \\
\hline
Tipos celulares L2 & 26 (25 biológicos + Unassigned) \\
\hline
ARI validación anotación & 0.742 \\
\hline
Genes con señal espacial (SVG consensus) & 288/289 (99.7\%) \\
\hline
Dominios espaciales (Novae) & 10 \\
\hline
Interacciones CCC significativas (LIANA+) & 1,565 \\
\hline
Interacciones validadas espacialmente (Spacia) & 22/25 (88\%) \\
\hline
Linajes pseudotemporales (Slingshot) & 8 \\
\hline
pval\_adj más significativo (Spacia ADAM17/MUC1) & $1.1 \times 10^{-71}$ \\
\hline
\end{tabular}

\vspace{1cm}
\begin{center}
\rule{0.5\textwidth}{0.4pt}\par
\vspace{0.3cm}
\textit{Fin del reporte técnico v3 --- 30/04/2026}
\end{center}

\end{document}
